%% S11: the converse of the tooth lemma against a run with no exposed move.
%% DRAFT 2026-11-25.  Neutral amsart shell, as S1--S9.
%% Source: ~/Code/ski/notes/LOCK1-CTB.md section 12 (the linear proof, Wizard-reviewed), with sections 0,
%%         11 (the set B, Lemmas CL, B, SOL) and 11.8 (where S1's block model enters); background
%%         ~/Code/ski/notes/NCFP-proof.md sections 0, 2 and LOCK1-CT.md sections 0-2.
%%         The corollary (rests of p_K) is the argument recorded in A1-lock1-S3.md, 2026-11-25, rechecked here.
%% Not mechanized.  Lean is cited only for the inputs (lean-ski v1.20.0, tag merge 20d8271):
%%   (CR) ....... church_rosser (Ski/Weak.lean)
%%   (Std) ...... wsteps_split (Ski/Weak.lean): Steps t t' and ParIs t' v; ParI keeps the spine
%%   (F2) ....... Full.run_total (Ski/Full.lean), through TopTerm.reach_all and Full.stable_iff (block 36 of S1)
%%   all three in #audit_clean (Ski/Audit.lean).  Lean's Term is SKI; S-terms are closed under its steps.
%% Map to the note (section 12): lem:head = Lemma L, (Std-h), (H), (AS); lem:moves = the Moves;
%%   lem:inv = S6 Lemma 4 (syntactic form) + Lemma T; lem:finite = Prop. 1 (<=); lem:suffix = Lemma R,
%%   RN-1, RN-2; lem:pk = Facts S, S'; lem:teeth = (K), (KT); lem:index = T1, T1a, T1b;
%%   lem:srun = Step 1; lem:cancel = Step 3 (SC); lem:left = Steps 2 and 4; lem:persist = Step 5, Lemma P;
%%   lem:closure = Lemma CL; lem:typing = Lemma B; lem:sol = Lemma SOL.
%% TODO before release:
%%   1. Wizard on this text; Grammarian on the prose.
%%   2. lean-ski v1.20.0: check the tag URL in refs.bib (leanski) and add the DOI.
%%   3. Recheck "Proposition 4, Theorem 9" of S1, "Theorem 7, Proposition 9" of S3 and "Lemma 4" of S6
%%      against the final versions.
\documentclass[11pt]{amsart}
\usepackage[T1]{fontenc}
\usepackage{lmodern}
\usepackage{amsmath,amssymb,amsthm}
\usepackage[numbers,sort]{natbib}
\usepackage[hidelinks]{hyperref}
\usepackage{bm}

\newtheorem{theorem}{Theorem}
\newtheorem{lemma}[theorem]{Lemma}
\newtheorem{corollary}[theorem]{Corollary}
\newtheorem{conjecture}[theorem]{Conjecture}
\theoremstyle{definition}
\newtheorem{definition}[theorem]{Definition}

\newcommand{\Ter}{\mathcal S}
\newcommand{\red}{\twoheadrightarrow}
\newcommand{\hred}{\twoheadrightarrow_h}
\newcommand{\conv}{=_\beta}
\newcommand{\sz}[1]{\lvert#1\rvert}
\newcommand{\rest}{\mathrm{rest}}
\newcommand{\pK}{p_K}
\newcommand{\B}{\mathfrak B}
\newcommand{\Red}{\mathrm{Red}}
\newcommand{\vU}{\vec U}

\title[Teeth Determine a Run That Exposes Nothing]{An $\bm{S}$-Term Determined by Its Teeth, Though No Tooth Exposes Its State}
\author{Claude Opus \& Fable with N.\ E.\ Davis}
\address{ZeroAttest, Champaign, Illinois, USA}
\email{NDavis@ZeroAttest.com}
\urladdr{https://orcid.org/0000-0002-1131-8319}
\date{\today}

\begin{document}
\begin{abstract}
The head run of a term built from the combinator $S$ that has no head normal form splits off an
infinite sequence of arguments, its teeth, and convertible terms have convertible teeth. Whether the
converse holds is open. We prove it against one run. Let $Q=S\,(S\,S)$, $G=S\,Q\,Q$ and
$p_K=S\,(S\,S)\,G\,(Q\,(Q\,G))\,S$. Every closed $S$-term without a head normal form whose $n$-th
tooth is convertible with that of $p_K$, for every $n$, is convertible with $p_K$; the same holds,
with the teeth shifted, for every state of $p_K$'s run. In this run no tooth exposes the state that produced it, so the
state cannot be read off a single tooth. Instead we show that a second run with the same teeth
must follow $p_K$'s spine equations, and that the reducts of $p_K$'s teeth lie among the reducts of
an explicit set of terms whose passive arguments are never put in function position; there the spine
equations have only $p_K$'s own solutions. The proof uses excluded middle in one form and is not mechanized.
\end{abstract}

\keywords{combinatory logic, combinator $S$, head normal form, convertibility, standardization}
\maketitle

\section{Introduction}

$\Ter$ is the set of closed terms built from $S$ by application, with $S\,x\,y\,z\to x\,z\,(y\,z)$ at
any position; $\red$ is the reflexive--transitive closure of $\to$, and $\conv$ is convertibility,
which by Church--Rosser is the existence of a common reduct. The \emph{size} $\sz t$ is the number of
occurrences of $S$ in $t$.

Every $t\in\Ter$ is $S\,u_1\cdots u_m$ for a unique $m\ge0$; the $u_i$ are its \emph{items}, and $t$
is a \emph{head normal form} (hnf) iff $m\le2$. If $m\ge3$, the \emph{head step} is
$t\to u_1\,u_3\,(u_2\,u_3)\,u_4\cdots u_m$. Head steps give the \emph{head run} of $t$, whose terms
are its \emph{frames}; $\hred$ is their reflexive--transitive closure. The term $t$ \emph{has an hnf} if its
head run reaches one, and then the first hnf reached is \emph{the} hnf of $t$; $t$ is \emph{literal} if
it is its own hnf. A frame $f\,w$ is \emph{stable} if $f$ has no hnf.

Let $t$ have no hnf. Its head run reaches a stable frame (input (F2) below); if $f\,w$ is the first,
put $\tau_1(t)=w$ and $\rest(t)=f$. Then $\rest(t)$ again has no hnf, and the \emph{teeth} of $t$ are
$\tau_{n+1}(t)=\tau_1(\rest^n(t))$; \emph{move} $n$ goes from $\rest^{n-1}(t)$ to $\rest^n(t)$ and
produces $\tau_n(t)$. If $t\conv t'$ have no hnf, then $\tau_n(t)\conv\tau_n(t')$ for all $n$
(Lemma~\ref{lem:inv}, the tooth lemma). Call \emph{CT} the converse: two terms without an hnf whose teeth are pairwise
convertible are convertible. Agreement of finitely many teeth says nothing about the states
$\rest^n$; but if the states ever meet, $\rest^n(p)\conv\rest^n(q)$, then $p\conv q$
(Lemma~\ref{lem:finite}). CT asks whether the tails force one meeting.

There is a simple case. At a move of kind (B) in Lemma~\ref{lem:moves} whose last hnf
$S\,a\,b$ is reached by at least one head step, the last frame is $S\,S\,y\,a$
(Lemma~\ref{lem:head}(a)), so the tooth is $(y\,a)\,w$ and the new state is $a\,w$. Call the move \emph{exposed} if $y$ has no hnf. Then $(y\,a)\,w$
and $y\,a$ are stable frames, and $\tau_1$ and $\rest$ read $w$ and $a$ off the tooth; if the
moves $n$ of $p$ and of $q$ are both exposed, their states are convertible and $p\conv q$. In the
run of the term $\pK$ below, $y$ is always a counter $Q^jG$ with a two-item hnf
(equation~\eqref{eq:am}), so no move is exposed, and nothing is read off a single tooth.

\begin{theorem}\label{thm:main}
Let $Q=S\,(S\,S)$, $G=S\,Q\,Q$ and $\pK=S\,(S\,S)\,G\,(Q\,(Q\,G))\,S$. Every $q\in\Ter$ without an
hnf such that $\tau_n(q)\conv\tau_n(\pK)$ for all $n\ge1$ satisfies $q\conv\pK$.
\end{theorem}

\begin{corollary}\label{cor:rests}
Let $n\ge0$. Every $q\in\Ter$ without an hnf such that $\tau_i(q)\conv\tau_{n+i}(\pK)$ for all
$i\ge1$ satisfies $q\conv\rest^n(\pK)$.
\end{corollary}

\emph{Assumptions.} Besides Church--Rosser and standardization, stated as (CR) and (Std) in
Section~\ref{sec:head}, the proof uses (F2): the head run of every term reaches an hnf or a stable
frame. (F2) is proved in~\citep{s1,s3} with the 38-element block model of~\citep{s1}, and this is the
only place where that model enters. The proof is classical in one respect: it splits cases on
whether a term has an hnf (in Lemma~\ref{lem:moves} and wherever a term's hnf is taken or its
absence assumed). The block model decides that property~\citep{s1}, so each such split could be made
constructive with it; we do not do so. Everything else is proved below. The proof is not
mechanized.

\emph{Plan.} After facts about head runs (Section~\ref{sec:head}) and the run of $\pK$
(Section~\ref{sec:pk}), Section~\ref{sec:q} puts a run with $\pK$'s teeth onto $\pK$'s spine
equations, and Section~\ref{sec:world} solves these equations inside the reducts of $\pK$'s teeth.

\section{Head runs}\label{sec:head}

A step is \emph{internal} if it contracts a redex inside an item. We use three known results.
\begin{itemize}
\item[(CR)] \citep{terese} If $t\red u$ and $t\red v$, then $u\red w$ and $v\red w$ for some $w$.
\item[(Std)] \citep{terese} If $t\red u$, then $t\hred t'\red u$ for some $t'$ such that every step
  from $t'$ to $u$ is internal. So $t'=S\,v_1\cdots v_m$ and $u=S\,v_1'\cdots v_m'$ with the same
  $m$ and $v_i\red v_i'$.
\item[(F2)] \citep[Proposition~9(b), $d=1$]{s3}, with \citep[Proposition~4 and Theorem~9]{s1}. The
  head run of every term reaches an hnf or a stable frame.
\end{itemize}
The head run is deterministic and stops at its first hnf. A right context of depth $d$ is
$x_1\,(x_2\,(\cdots(x_d\,[\,])\cdots))$; $e$ is a \emph{right suffix of $t$ at depth $d$} if $t$ is
such a context filled with $e$.

\begin{lemma}\label{lem:head}
\textup{(a)} A head step from $S\,u_1\cdots u_m$ gives a frame whose items are those of $u_1$,
followed by $u_3,\ u_2\,u_3,\ u_4,\dots,u_m$. Hence if a head run takes at least one step to a
two-item hnf, its last frame is $S\,S\,y\,z$ and the hnf is $S\,z\,(y\,z)$.

\textup{(b)} While the frames of $f$ have at least three items, the head run of $f\,w$ is that of $f$
with $w$ appended. A term without an hnf has only frames with at least three items.

\textup{(c)} A term has an hnf iff it reduces to an hnf. A term that head-reduces to a term without
an hnf has none.

\textup{(d)} If $f$ has no hnf, then neither has $f\,w$, nor any reduct of $f$.

\textup{(e)} Every reduct of an application is an application. Hence $t\conv S$ implies $t=S$.
\end{lemma}

\begin{proof}
(a) Two items need $m=3$ and $u_1=S$. (b) A head step from at least three items acts on the
prefix; a frame with at most two items is an hnf. (c) If $t\red h$ with $h$ an hnf, (Std) gives
$t\hred t'\red h$ with internal steps after $t'$, which keep the number of items, so $t'$ is an hnf.
For the second claim, if $t\hred t'$ and the run of $t$ reaches an hnf, then $t'$ lies on the run
before it, so $t'$ has an hnf.
(d) By (b) the run of $f\,w$ is that of $f$ with $w$ appended and reaches no hnf; if a reduct of
$f$ had an hnf, $f$ would reduce to an hnf, and (c) applies. (e) A step at the root of
$S\,x\,y\,z$ gives the application $x\,z\,(y\,z)$, and other steps keep the root. A common reduct of
$t$ and the normal form $S$ is $S$, so $t\red S$, and $t$ is not an application.
\end{proof}

\begin{lemma}\label{lem:moves}
Let $f=f^-\,u$ have no hnf.

\textup{(A)} If $f^-$ has no hnf, then $f$ is stable, $\tau_1(f)=u$ and $\rest(f)=f^-$.

\textup{(B)} Otherwise $f^-\hred S\,a_0\,b_0$, and with $w_0=u$, $w_{i+1}=b_i\,w_i$ there is $k\ge0$
such that $a_i\,w_i\hred S\,a_{i+1}\,b_{i+1}$ for $i<k$, $a_k\,w_k$ has no hnf, $\tau_1(f)=b_k\,w_k$
and $\rest(f)=a_k\,w_k$. So $w_k,w_{k-1},\dots,w_0=u$ are right suffixes of $\tau_1(f)$ at depths
$1,2,\dots,k+1$, and $w_k$ is an application if $k\ge1$.
\end{lemma}

\begin{proof}
(A) $f$ is a frame with at least three items whose prefix has no hnf. (B) By
Lemma~\ref{lem:head}(b) the run of $f$ follows that of $f^-$, with $u$ appended, to the hnf $h$ of
$f^-$; if $h$ had at most one item, $h\,u$ would be an hnf. So $f\hred S\,a_0\,b_0\,w_0\to
a_0\,w_0\,w_1$. If $a_0\,w_0$ has no hnf, this frame is stable. Otherwise the run follows that of
$a_0\,w_0$ to its hnf, which has two items for the same reason, and reaches
$S\,a_1\,b_1\,w_1$; and so on. Every frame passed before the construction stops is $c\,w_j$ with $c$
a frame of a term that has an hnf, or $S\,a_i\,b_i\,w_i$, so none is stable. By (F2) the
construction stops.
\end{proof}

\begin{lemma}\label{lem:inv}
Let $t\conv t'$. \textup{(a)} If $t$ has the hnf $S\,u_1\cdots u_m$, then $t'$ has an hnf
$S\,u'_1\cdots u'_m$, with the same $m$ and $u_i\conv u'_i$. \textup{(b)} If $t$ has no hnf, then
$t'$ has none, $\tau_1(t)\conv\tau_1(t')$ and $\rest(t)\conv\rest(t')$; hence
$\tau_n(t)\conv\tau_n(t')$ and $\rest^n(t)\conv\rest^n(t')$ for all $n$.
\end{lemma}

\noindent This is, in substance, \citep[Lemma~4]{s6} with its proof; we prove it from (CR), (Std)
and (F2) alone.

\begin{proof}
(a) Every reduct of $T=S\,u_1\cdots u_m$ ($m\le2$) is $S\,u^*_1\cdots u^*_m$ with $u_i\red u_i^*$.
By (CR), $t'$ and $T$ have a common reduct $T^*$; by (Std), $t'\hred T'\red T^*$ with internal steps,
so $T'$ has $m$ items, each reducing to the corresponding item of $T^*$, and $T'$ is the hnf of $t'$.
(b) By (a), $t'$ has no hnf. Let $T=f\,w$ be the first stable frame of $t$. A reduct of $T$ is
$f^*\,w^*$ with $f\red f^*$, $w\red w^*$, since $f^*$ has no hnf (Lemma~\ref{lem:head}(d)) and so is not
$S\,a\,b$. As in (a), $t'\hred T'=f'\,w'\red f^*\,w^*$ with $f'\red f^*$, $w'\red w^*$. Now
$f'\conv f$ has no hnf by (a), so $T'$ is stable. The first stable frame $f''\,w''$ of $t'$ is at
or before $T'$, and by Lemma~\ref{lem:head}(b) every later head step acts on the prefix; so $w''=w'$
and $f''\hred f'$. Hence $w''\conv w$ and $f''\conv f$. Induction gives the rest, since $\rest(t)$ has
no hnf and $\tau_{n+1}(t)=\tau_1(\rest^n(t))$.
\end{proof}

\begin{lemma}\label{lem:finite}
Let $p,q$ have no hnf. If $\rest^n(p)\conv\rest^n(q)$ and $\tau_i(p)\conv\tau_i(q)$ for $i\le n$,
then $p\conv q$.
\end{lemma}

\begin{proof}
$p\hred\rest(p)\,\tau_1(p)$, and $\rest(p)$ has no hnf, so by Lemma~\ref{lem:head}(b) its run
continues with $\tau_1(p)$ appended. By induction $p\hred\rest^n(p)\,\tau_n(p)\cdots\tau_1(p)$, likewise
for $q$, and $\conv$ is a congruence.
\end{proof}

\begin{lemma}\label{lem:suffix}
\textup{(a)} If $t\red t'$ and $e$ is a right suffix of $t$ at depth $d$, then $t'$ has a right
suffix $e'$ at a depth $\ge d$ with $e\red e'$.

\textup{(b)} Let $t_0=t$ and $t_{i+1}=\tau_1(t_i)$, where $t_0,\dots,t_{J-1}$ have no hnf and
$t_J=S$ \textup{(}$J\ge1$\textup{)}. Every right suffix of $t$ at depth $d$ is $S$ or is convertible
with $t_i$ for some $i$ with $d\le i<J$.

\textup{(c)} If $t\conv Z$, where $Z$ has no hnf and $\tau_1(Z)=S$, then every right suffix of $t$
at depth $\ge1$ is $S$.

\textup{(d)} If $t\conv Y$, where $Y$ and $\tau_1(Y)$ have no hnf and $\tau_1(\tau_1(Y))=S$, then
every right suffix of $t$ at depth $1$ is $S$ or convertible with $\tau_1(Y)$, and every one at
depth $\ge2$ is $S$.
\end{lemma}

\begin{proof}
(a) One step: if the redex lies inside $e$, take its contractum in place; if it lies off the path
to $e$, $e$ stays at its depth; otherwise the redex is a right suffix $S\,a\,b\,c$ with $e$ inside
$c$, and in $a\,c\,(b\,c)$ the term $c$, hence $e$, moves one level down.
(b) Let $f_i\,t_{i+1}$ be the first stable frame of $t_i$ ($i<J$). Put $r_J=S$ and
$r_i=f_i\,r_{i+1}$; by downward induction $t_i\red r_i$, so $t\red r_0$, whose right suffixes are
$r_0,\dots,r_J$ at depths $0,\dots,J$. By (a), $e\red r_{i}$ for some $i\ge d$. If $i=J$, then $e=S$
(Lemma~\ref{lem:head}(e)); otherwise $e\conv r_i\conv t_i$.
(c) By Lemma~\ref{lem:inv}, $t$ has no hnf and $\tau_1(t)\conv S$, so $\tau_1(t)=S$; apply (b) with
$J=1$. (d) Likewise $\tau_1(t)\conv\tau_1(Y)$ has no hnf and $\tau_1(\tau_1(t))=S$; apply (b) with
$J=2$.
\end{proof}

\section{The run of \texorpdfstring{$p_K$}{pK}}\label{sec:pk}

Let $c_0=G$ and $c_{j+1}=Q\,c_j$, so $\pK=Q\,c_0\,c_2\,S$. For all $u,v\in\Ter$ the head runs
\begin{equation}\label{eq:h}
\begin{gathered}
Q\,u\,v\to S\,S\,v\,(u\,v)\to S\,(u\,v)\,(v\,(u\,v)),\\
G\,u\to Q\,u\,(Q\,u)\hred S\,(u\,(Q\,u))\,(Q\,u\,(u\,(Q\,u)))
\end{gathered}
\end{equation}
are literal. Put $a_0=Q\,c_0\,c_2$. If $a_m=Q\,c_j\,c_n$, put $a_{m+1}=c_j\,c_n$ and
$b_{m+1}=c_n\,a_{m+1}$; if $a_m=G\,c_n$, put $a_{m+1}=c_n\,c_{n+1}$ and
$b_{m+1}=c_{n+1}\,a_{m+1}$. Since $c_j\,c_n$ is $Q\,c_{j-1}\,c_n$ or $G\,c_n$, and $c_n\,c_{n+1}=Q\,c_{n-1}\,c_{n+1}$,
every $a_m$ is $Q\,c_j\,c_n$ or $G\,c_n$ with $n\ge2$, and by~\eqref{eq:h}
\begin{equation}\label{eq:am}
a_m\hred S\,S\,c\,a_{m+1}\to S\,a_{m+1}\,b_{m+1},\qquad c\in\{c_n,c_{n+1}\}.
\end{equation}
So $b_m=c_{i_m+1}\,a_m$ with $i_m\ge1$. The index $n$ never decreases and grows at every step from
$G\,c_n$; such steps recur, since $j$ falls at every step from $Q\,c_j\,c_n$; and $i_{m+1}\ge n-1$.
So \emph{the sequence $i_m$ is unbounded}. Put $x_m=c_{i_m}\,a_m$ and $Y_n=b_n\,S$. By~\eqref{eq:h},
\begin{equation}\label{eq:bm}
b_m=Q\,c_{i_m}\,a_m\hred S\,x_m\,(a_m\,x_m).
\end{equation}

\begin{lemma}\label{lem:pk}
\textup{(a)} For all $m$ and $u\in\Ter$, $a_m\,u$ has no hnf, and its first stable frame is
$a_{m+1}\,u\,(b_{m+1}\,u)$.

\textup{(b)} For $r\ge1$ and $U_1,\dots,U_r\in\Ter$, the terms $a_m\,\vU$, $c_j\,a_m\,\vU$
\textup{(}$j\ge0$\textup{)} and $Q\,a_m\,(a_m\,(Q\,a_m))\,\vU$ have no hnf, where
$\vU=U_1\cdots U_r$.
\end{lemma}

\begin{proof}
(a) By~\eqref{eq:am} and Lemma~\ref{lem:head}(b), $a_m\,u\hred a_{m+1}\,u\,(b_{m+1}\,u)$ in at least
one step, through frames with at least three items. Repeating with appended arguments,
$a_m\,u\hred a_{m+j}\,u\,V_j$ in at least $j$ steps for every $j$, so the head run is infinite and
reaches no hnf. Before $a_{m+1}\,u\,(b_{m+1}\,u)$, every frame is $c\,u$ with $c$ a frame of $a_m$, or
$S\,a_{m+1}\,b_{m+1}\,u$; none is stable. (b) The first by (a) and Lemma~\ref{lem:head}(d). By~\eqref{eq:h},
$c_{j}\,a\,\vU\hred c_{j-1}\,a\,U_1\,(a\,(c_{j-1}\,a)\,U_1)\,U_2\cdots U_r$ for $j\ge1$,
$G\,a\,\vU\hred a\,(Q\,a)\,U_1\cdots$ and $Q\,a\,v\,\vU\hred a\,v\,U_1\cdots$; so each head-reduces to
a term of the first kind, and Lemma~\ref{lem:head}(c) applies.
\end{proof}

\begin{lemma}\label{lem:teeth}
For $n\ge1$: $\rest^{n}(\pK)=a_{n}\,S$ and $\tau_n(\pK)=Y_n$. The terms $Y_n$, $x_n\,S$, $a_n\,x_n$
and $a_n\,x_n\,S$ have no hnf, $\rest(Y_n)=x_n\,S$, $\tau_1(Y_n)=a_n\,x_n\,S$,
$\rest(\tau_1(Y_n))=a_n\,x_n$ and $\tau_1(\tau_1(Y_n))=S$. So Lemma~\ref{lem:suffix}\textup{(d)}
applies to every $t\conv Y_n$, and Lemma~\ref{lem:suffix}\textup{(c)} to every $t\conv\tau_1(Y_n)$.
\end{lemma}

\begin{proof}
The first sentence is Lemma~\ref{lem:pk}(a) with $u=S$. The terms have no hnf by
Lemma~\ref{lem:pk}(b), since $Y_n=c_{i_n+1}\,a_n\,S$. By~\eqref{eq:bm} and Lemma~\ref{lem:moves}(B),
the move on $Y_n$ has $a_0\,w_0=x_n\,S$, without an hnf, so $k=0$; the move on $a_n\,x_n\,S$ is
Lemma~\ref{lem:moves}(A).
\end{proof}

\noindent By~\eqref{eq:am}, the last frame of move $n$ of $\pK$ is $S\,S\,c\,a_n$ with
$c\in\{c_n,c_{n+1}\}$, which has a two-item hnf: no move is exposed.

For $x$ with an hnf $S\,u_1\cdots u_m$, $m\ge1$, put $x\cdot1=u_1$, and let $\kappa(x)$ be the least
$k$ such that $x\cdot1^k$ is defined and has no hnf.

\begin{lemma}\label{lem:index}
\textup{(a)} If $x\conv x'$, then $\kappa(x)=\kappa(x')$, both sides being defined or not.
\textup{(b)} $\kappa(x_n)=i_n+1$. \textup{(c)} If $\tau_1(Y_n)\conv\tau_1(Y_{n'})$, then $i_n=i_{n'}$.
\end{lemma}

\begin{proof}
(a) Lemma~\ref{lem:inv}(a), by induction along $1^k$. (b) By~\eqref{eq:h}, we have
$c_j\,a\hred S\,(c_{j-1}\,a)\cdots$ for $j\ge1$, and $c_0\,a=G\,a\hred S\,(a\,(Q\,a))\cdots$, where
$a\,(Q\,a)$ has no hnf (Lemma~\ref{lem:pk}(b)). (c) By Lemma~\ref{lem:teeth},
$\rest(\tau_1(Y_n))=a_n\,x_n$, and by Lemma~\ref{lem:pk}(a), $\tau_1(a_n\,x_n)=b_{n+1}\,x_n$.
By~\eqref{eq:bm} and Lemma~\ref{lem:moves}(B), where $a_0\,w_0=x_{n+1}\,x_n$ has no hnf
(Lemma~\ref{lem:pk}(b)),
$\tau_1(b_{n+1}\,x_n)=a_{n+1}\,x_{n+1}\,x_n$, and by Lemma~\ref{lem:moves}(A) its $\tau_1$ is $x_n$.
All these terms have no hnf, so Lemma~\ref{lem:inv}(b), applied four times, gives
$x_n\conv x_{n'}$, and (a), (b) give $i_n=i_{n'}$.
\end{proof}

\section{Runs with the teeth of \texorpdfstring{$p_K$}{pK}}\label{sec:q}

From now on, $q\in\Ter$ has no hnf and $\tau_n(q)\conv Y_n$ for all $n\ge1$. By
Lemma~\ref{lem:inv}(b), no $\tau_n(q)$ has an hnf. Write $f_n=\rest^n(q)$.

\begin{lemma}\label{lem:srun}
There is $n_1\ge1$ such that, for every $n\ge n_1$, $f_{n-1}=a'_{n-1}\,S$ with
$a'_{n-1}\hred S\,a'_n\,b'_n$, $\tau_n(q)=b'_n\,S$ and $f_n=a'_n\,S$.
\end{lemma}

\begin{proof}
Consider a move $n$ of $q$ of kind (B) in Lemma~\ref{lem:moves}, with $w_k$ its last $w$. By
Lemma~\ref{lem:moves}, $w_k$ is a right suffix of $\tau_n(q)\conv Y_n$ at depth $1$, so by
Lemmas~\ref{lem:teeth} and~\ref{lem:suffix}(d) it is $S$ or $w_k\conv\tau_1(Y_n)$.

The second case never occurs. Suppose it does, put $w=w_k$, so $f_n=a_k\,w$. If move $n+1$ were of
kind (A), then $\tau_{n+1}(q)=w$, and Lemma~\ref{lem:inv}(b) would give
$\tau_1(w)\conv\tau_1(\tau_1(Y_n))=S$ and $\tau_1(w)\conv\tau_1(Y_{n+1})$, which has no hnf.
So move $n+1$ is of kind (B) with $u=w$, and $w$ is a right suffix of $\tau_{n+1}(q)\conv Y_{n+1}$ at
depth $k'+1$, where $k'$ is the $k$ of that move. As $w\ne S$, Lemma~\ref{lem:suffix}(d) forces
$k'=0$ and $w\conv\tau_1(Y_{n+1})$, and $w$ is again the last $w$ of move $n+1$. By induction
$\tau_1(Y_{n'})\conv\tau_1(Y_n)$ for all $n'\ge n$, so $i_{n'}=i_n$ by Lemma~\ref{lem:index}(c),
against the unboundedness of $i$.

So at every move of kind (B), $w_k=S$; as $w_k$ is an application if $k\ge1$, $k=0$ and $u=S$. A
move of kind (A) makes the state a proper subterm, so some move $n_1$ is of kind (B). After a
move of kind (B) the state is $a_0\,S$; a move of kind (A) from it would make $S$ a tooth, which has an
hnf. So every move $n\ge n_1$ is of kind (B) with $k=0$ and $u=S$, which is the claim.
\end{proof}

\begin{lemma}\label{lem:cancel}
\textup{(a)} If $t$ has no hnf and $e\,S\conv t\,S$, then $e\conv t$.

\textup{(b)} Let $h\hred S\,u\,v$, where $u\,S$ has no hnf and $u$ or $v$ has no hnf. If
$e\,S\conv h\,S$, then $e\hred S\,e_1\,e_2$ with $e_1\,S\conv u\,S$ and $e_2\,S\conv v\,S$.
\end{lemma}

\begin{proof}
We use that $\tau_1$ of a move of kind (B) is an application (Lemma~\ref{lem:moves}).
(a) By Lemma~\ref{lem:moves}(A), $\tau_1(t\,S)=S$ and $\rest(t\,S)=t$. By Lemma~\ref{lem:inv},
$e\,S$ has no hnf, so $e$ has no hnf with at most one item; if $e$ had an hnf with two items, $\tau_1(e\,S)$
would be an application convertible with $S$, against Lemma~\ref{lem:head}(e). So $e$ has no hnf, and
$e=\rest(e\,S)\conv t$.

(b) By Lemma~\ref{lem:moves}(B) with $k=0$, $\tau_1(h\,S)=v\,S$ and $\rest(h\,S)=u\,S$. As in (a),
$e$ has no hnf with at most one item; and $e$ has an hnf, since otherwise $\tau_1(e\,S)=S$ would be
convertible with the application $v\,S$. So $e\hred S\,e_1\,e_2$, and the move on $e\,S$
has $a_0=e_1$, $b_0=e_2$, $w_0=S$. If its $k$ is $0$, Lemma~\ref{lem:inv}(b) gives the claim. Let
$k\ge1$. If $v$ has no hnf, then $v\,S$ has $\tau_1(v\,S)=S$, and the application $w_1=e_2\,S$ is a
right suffix of $\tau_1(e\,S)\conv v\,S$ at depth $k\ge1$, against Lemma~\ref{lem:suffix}(c). If $u$
has no hnf, then $\rest(e\,S)=a_k\,w_k\conv u\,S$, whose $\tau_1$ is $S$, so
$\tau_1(a_k\,w_k)\conv S$ by Lemma~\ref{lem:inv}(b); but $\tau_1(a_k\,w_k)$ is $w_k$ or the $\tau_1$ of a
move of kind (B), an application in either case.
\end{proof}

\begin{lemma}\label{lem:left}
For every $n\ge n_1$, $b'_n\conv b_n$.
\end{lemma}

\begin{proof}
Write $b'=b'_n$, $a=a_n$ and $i=i_n$, so $x_n=c_i\,a$. By Lemma~\ref{lem:inv}(b),
$b'\,S\conv Y_n$ gives $\tau_1(b'\,S)\conv a\,x_n\,S$, which has no hnf. If $b'$ had no hnf,
$\tau_1(b'\,S)$ would be $S$, which has an hnf; with an hnf of at most one item, $b'\,S$ would be an hnf. So
$b'\hred S\,\alpha\,\beta$, and the move on $b'\,S$ has $a_0=\alpha$, $b_0=\beta$, $w_0=S$. If its $k$
were $\ge1$, then $w_k$, an application, would be a right suffix at depth $1$ of $\tau_1(b'\,S)$,
against Lemmas~\ref{lem:teeth} and~\ref{lem:suffix}(c). So $k=0$, and Lemma~\ref{lem:inv}(b) gives
$\alpha\,S\conv x_n\,S$ and $\beta\,S\conv a\,x_n\,S$. By Lemma~\ref{lem:cancel}(a), with $t=a\,x_n$,
$\beta\conv a\,x_n$.

It remains to show $\alpha\conv x_n$. Put $e_0=\alpha$, so $e_0\,S\conv c_i\,a\,S$. For $k<i$,
suppose $e_k\,S\conv c_{i-k}\,a\,S$. By~\eqref{eq:h},
$c_{i-k}\,a\hred S\,(c_{i-k-1}\,a)\,(a\,(c_{i-k-1}\,a))$, and $c_{i-k-1}\,a\,S$ and
$a\,(c_{i-k-1}\,a)$ have no hnf (Lemma~\ref{lem:pk}(b)). So Lemma~\ref{lem:cancel}(b) gives
$e_k\hred S\,e_{k+1}\,e'_{k+1}$ with $e_{k+1}\,S\conv c_{i-k-1}\,a\,S$ and, by
Lemma~\ref{lem:cancel}(a), $e'_{k+1}\conv a\,(c_{i-k-1}\,a)$. At $k=i$ we have $e_i\,S\conv G\,a\,S$,
and $G\,a\hred S\,v\,(Q\,a\,v)$ with $v=a\,(Q\,a)$; here $v$ and $v\,S$ have no hnf. So
Lemma~\ref{lem:cancel}(b), then (a), give $e_i\hred S\,\zeta\,\eta$ with $\zeta\conv v$ and
$\eta\,S\conv Q\,a\,v\,S$. Finally $Q\,a\,v\hred S\,(a\,v)\,(v\,(a\,v))$, and $a\,v$, $a\,v\,S$ and
$v\,(a\,v)$ have no hnf; so, in the same way, $\eta\hred S\,\eta_1\,\eta_2$ with $\eta_1\conv a\,v$ and
$\eta_2\conv v\,(a\,v)$. Reassembling,
$\eta\conv S\,(a\,v)\,(v\,(a\,v))\conv Q\,a\,v$, then $e_i\conv S\,v\,(Q\,a\,v)\conv G\,a$, then $e_k\conv c_{i-k}\,a$ for
$k=i-1,\dots,0$. So $\alpha\conv x_n$, and
$b'\conv S\,\alpha\,\beta\conv S\,x_n\,(a\,x_n)\conv b_n$ by~\eqref{eq:bm}.
\end{proof}

\section{The world of \texorpdfstring{$p_K$}{pK}}\label{sec:world}

\begin{definition}
A \emph{solution from $k$} is a term $s=s_0$ with $s_j\hred S\,s_{j+1}\,t_{j+1}$ (its hnf) and
$t_{j+1}\conv b_{k+j+1}$ for all $j\ge0$.
\end{definition}

\begin{lemma}\label{lem:persist}
\textup{(a)} A term convertible with a solution from $k$ is one.
\textup{(b)} $a_k$ is a solution from $k$, with $s_j=a_{k+j}$ and $t_j=b_{k+j}$; and $a'_k$ is a
solution from $k$ for every $k\ge n_1$.
\textup{(c)} Let $s$ be a solution. If $t$ contains a reduct of $s_j$, then every reduct of $t$
contains a reduct of some $s_{j'}$ with $j'\ge j$.
\end{lemma}

\begin{proof}
(a) Lemma~\ref{lem:inv}(a), by induction on $j$. (b) By~\eqref{eq:am}, and by
Lemmas~\ref{lem:srun} and~\ref{lem:left}. (c) Take one step $t\to t'$ that contracts
$\rho=S\,x\,y\,z$, and an occurrence $o$ in $t$ of a reduct $r$ of $s_j$. If $o$ is disjoint from
$\rho$, lies in $x$, $y$ or $z$, or contains $\rho$, then $t'$ contains $o$ or its one-step reduct.
Otherwise $o$ is $S$, $S\,x$ or $S\,x\,y$. The first two are impossible: $r\conv s_j$ has a
two-item hnf (Lemma~\ref{lem:inv}(a)). If $r=S\,x\,y$, then (Std) gives
$s_j\hred S\,s_{j+1}\,t_{j+1}\red S\,x\,y$ by internal steps, so $x$ is a reduct of $s_{j+1}$, and
$x$ occurs in the contractum $x\,z\,(y\,z)$. Induction on the number of steps gives the claim.
\end{proof}

\emph{The set $\B$.} Let $U,U_1,\dots,U_r$ range over $\B$, with $r\ge0$ and $\vU=U_1\cdots U_r$. Then
$\B$ is the least set of terms that contains the \emph{pieces} $S$, $S\,S$, $Q$, $S\,Q$ and $c_j$
($j\ge0$), and the terms
\begin{center}\small$
\begin{array}{llll}
(\mathrm{Fa})& a_m\,\vU, &(\mathrm{Fx})& c_j\,a_m\,\vU,\\
(\mathrm{FQa})& Q\,a_m, &(\mathrm{FQa2})& Q\,a_m\,(a_m\,(Q\,a_m))\,\vU,\\
(\mathrm{FSS})& \multicolumn{3}{l}{S\,S\,w,\quad w\in\{c_j,\ a_m,\ Q\,a_m,\ a_m\,(Q\,a_m)\},}\\
(\mathrm{FS})& \multicolumn{3}{l}{S\,w,\quad w\in\{a_m,\ c_j\,a_m,\ a_m\,(Q\,a_m),\ a_m\,(a_m\,(Q\,a_m))\},}
\end{array}$
\end{center}
for all $j,m\ge0$. The $U_i$ are the \emph{passive} arguments. Each $b_n=c_{i_n+1}\,a_n$ is in
(Fx). Let $\Red(\B)$ be the set of reducts of members of $\B$.

\begin{lemma}\label{lem:closure}
Suppose that \textup{(i)} the function part and the argument of every member of $\B$ that is an
application are in $\B$, and \textup{(ii)} for every $\beta\in\B$ and every frame
$S\,x_0\,y_0\,z_0$ with exactly three items on the head run of $\beta$, the terms $x_0\,z_0$ and
$y_0\,z_0$ are in $\B$. Then $\Red(\B)$ is closed under subterms.
\end{lemma}

\begin{proof}
We show by induction on $n$ that every subterm $t$ of an $n$-step reduct $r$ of a member of $\B$ is
in $\Red(\B)$. For $n=0$ use (i). Otherwise $r_0\to r$ contracts $S\,x\,y\,z$ and $r_0$ is an
$(n-1)$-step reduct. If $t$ is disjoint from the contractum or lies in a copy of $x$, $y$ or $z$, it is
a subterm of $r_0$. If $t$ contains the contractum, it is a reduct of a subterm of $r_0$. Otherwise
$t$ is $x\,z$ or $y\,z$. By induction $S\,x\,y\,z\in\Red(\B)$, so $\beta\red S\,x\,y\,z$ for some
$\beta\in\B$, and (Std) gives $\beta\hred S\,x_0\,y_0\,z_0\red S\,x\,y\,z$ by internal steps, with
$x_0\red x$, $y_0\red y$, $z_0\red z$. So $x\,z$ is a reduct of $x_0\,z_0\in\B$, and likewise $y\,z$.
\end{proof}

\begin{lemma}\label{lem:typing}
$\B$ satisfies \textup{(i)} and \textup{(ii)} of Lemma~\ref{lem:closure}.
\end{lemma}

\begin{proof}
(i) Inspection: $c_0=S\,Q\,Q$, $c_{j+1}=Q\,c_j$, $Q=S\,(S\,S)$; $a_m$ is $c_{j+1}\,c_n$ or $c_0\,c_n$;
the function parts of members with $r\ge1$ are members of the same family with $r-1$; and
$c_j\,a_m$, $Q\,a_m$, $Q\,a_m\,(a_m\,(Q\,a_m))$, $S\,S\,w$ and $S\,w$ split into pieces and members of
(Fa), (Fx), (FQa).

(ii) Pieces and members of (FQa), (FSS) and (FS) have at most two items. For the others the head run
follows from~\eqref{eq:h} and Lemma~\ref{lem:head}(b); we write $a=a_m$, $a'=a_{m+1}$, $b'=b_{m+1}$,
$v=a\,(Q\,a)$, and the head of a member is the member with $r=0$. The run first passes the frames of the head
with $\vU$ appended, which have at least three items, exactly three only if $r=0$, until the head's
hnf $S\,g_1\,g_2$; then $S\,g_1\,g_2\,\vU$, with three items exactly if $r=1$; then, if $r\ge1$,
$g_1\,U_1\,(g_2\,U_1)\,U_2\cdots U_r$, which is again of the form (Fa) or (Fx), with $r+1\ge2$
arguments. Applying the same description to it, and so on, no later frame has exactly three items. So the frames with exactly three items, and
their terms $x_0\,z_0$; $y_0\,z_0$, are these (members written without $U$ have $r=0$):
\begin{center}\small$
\begin{array}{@{}lll@{}}
\text{member} & \text{frame } S\,x_0\,y_0\,z_0 & x_0\,z_0;\ y_0\,z_0\\[2pt]
a=Q\,c_j\,c_n & Q\,c_j\,c_n,\quad S\,S\,c_n\,a' & S\,S\,c_n;\ a',\quad S\,a';\ b'\\
a=G\,c_n & G\,c_n,\quad Q\,c_n\,c_{n+1} & c_{n+1};\ c_{n+1},\quad S\,S\,c_{n+1};\ a'\\
 & S\,S\,c_{n+1}\,a' & S\,a';\ b'\\
a\,U & S\,a'\,b'\,U & a'\,U;\ b'\,U\\
c_j\,a,\ j\ge1 & Q\,c_{j-1}\,a,\quad S\,S\,a\,(c_{j-1}\,a) & S\,S\,a;\ c_{j-1}\,a,\quad S\,(c_{j-1}\,a);\ a\,(c_{j-1}\,a)\\
c_j\,a\,U,\ j\ge1 & S\,(c_{j-1}\,a)\,(a\,(c_{j-1}\,a))\,U & c_{j-1}\,a\,U;\ a\,(c_{j-1}\,a)\,U\\
G\,a & G\,a,\quad Q\,a\,(Q\,a) & Q\,a;\ Q\,a,\quad S\,S\,(Q\,a);\ v\\
 & S\,S\,(Q\,a)\,v & S\,v;\ Q\,a\,v\\
G\,a\,U & S\,v\,(Q\,a\,v)\,U & v\,U;\ Q\,a\,v\,U\\
Q\,a\,v & Q\,a\,v,\quad S\,S\,v\,(a\,v) & S\,S\,v;\ a\,v,\quad S\,(a\,v);\ v\,(a\,v)\\
Q\,a\,v\,U & S\,(a\,v)\,(v\,(a\,v))\,U & a\,v\,U;\ v\,(a\,v)\,U
\end{array}$
\end{center}
Every term on the right is in $\B$: $c_{n+1}$ is a piece; $S\,S\,c_n$, $S\,S\,c_{n+1}$, $S\,S\,a$,
$S\,S\,(Q\,a)$ and $S\,S\,v$ are in (FSS); $S\,a'$, $S\,(c_{j-1}\,a)$, $S\,v$ and $S\,(a\,v)$ in (FS);
$Q\,a$ in (FQa); $Q\,a\,v$ and $Q\,a\,v\,U$ in (FQa2); $b'=c\,a'$, $c_{j-1}\,a$ and $c_{j-1}\,a\,U$ in
(Fx); and $a'$, $a'\,U$, $a\,(c_{j-1}\,a)$, $a\,(c_{j-1}\,a)\,U$, $v$, $v\,U$, $a\,v$, $a\,v\,U$,
$v\,(a\,v)$ and $v\,(a\,v)\,U$ in (Fa).
\end{proof}

\noindent In these runs a passive argument is only ever passed on as an argument ($z_0$ above).

\begin{lemma}\label{lem:sol}
If $\beta\in\B$ is a solution from $k$, then $\beta=a_k$.
\end{lemma}

\begin{proof}
A solution has a two-item hnf whose first item has a two-item hnf, and its second item $t_1\conv b_{k+1}$
has an hnf, by~\eqref{eq:bm} and Lemma~\ref{lem:inv}(a). Members of (Fa), (Fx) and (FQa2) with
$r\ge1$ have no hnf (Lemma~\ref{lem:pk}(b)). $S$, $S\,S$, $S\,Q$, $Q$ and the members of (FS) have at
most one item. The first item of $c_j$ ($j\ge1$), $G$, $Q\,a$ and $S\,S\,w$ is $S\,S$, $Q$ or $S$, with at
most one item. By~\eqref{eq:h}, the hnfs of $Q\,a\,v$ and $G\,a$ have the first items $a\,v$ and $a\,(Q\,a)$, which
have no hnf, and that of $c_j\,a$ ($j\ge1$) has the second item $a\,(c_{j-1}\,a)$, which has no hnf. So
$\beta=a_m$ for some $m$, with $t_j=b_{m+j}$; hence $b_{m+j}\conv b_{k+j}$ for all $j\ge1$.
By~\eqref{eq:bm} and Lemma~\ref{lem:inv}(a), $x_{m+j}\conv x_{k+j}$, and by
Lemma~\ref{lem:index}(a),(b), $i_{m+j}=i_{k+j}$ for all $j\ge1$. If $m\ne k$, the sequence $i$ would
be periodic from some index on, hence bounded. So $m=k$.
\end{proof}

\begin{proof}[Proof of Theorem~\ref{thm:main}]
If $a'_{n-1}$ is literal, that is, $a'_{n-1}=S\,a'_n\,b'_n$, then $\sz{a'_n}<\sz{a'_{n-1}}$. So
$a'_{n-1}$ is not literal for some $n>n_1$, and then Lemma~\ref{lem:head}(a) gives $b'_n=y\,a'_n$ for some $y$. By
Lemma~\ref{lem:left} and (CR), $b'_n\red R$ and $b_n\red R$ for some $R$. The term $b'_n$ contains
$a'_n$, which is $s_{n-n_1}$ for the solution $s=a'_{n_1}$ (Lemma~\ref{lem:persist}(b)). By
Lemma~\ref{lem:persist}(c), $R$ contains a reduct $s'$ of some $a'_k$ with $k\ge n$, and $s'$ is a
solution from $k$ by Lemma~\ref{lem:persist}(a),(b). As $b_n\in\B$, Lemmas~\ref{lem:closure}
and~\ref{lem:typing} give $s'\in\Red(\B)$: $\beta\red s'$ for some $\beta\in\B$. Then $\beta$ is a
solution from $k$ (Lemma~\ref{lem:persist}(a)), so $\beta=a_k$ by Lemma~\ref{lem:sol}, and
$a'_k\conv s'\conv a_k$. Hence $\rest^k(q)=a'_k\,S\conv a_k\,S=\rest^k(\pK)$ (Lemmas~\ref{lem:srun}
and~\ref{lem:teeth}), and Lemma~\ref{lem:finite} gives $q\conv\pK$.
\end{proof}

\begin{proof}[Proof of Corollary~\ref{cor:rests}]
Put $q'=q\,Y_n\cdots Y_1$. By Lemma~\ref{lem:head}(d), $q\,Y_n\cdots Y_{j+1}$ has no hnf for every
$j\le n$, so for $j<n$ it is a stable frame, its own first one; hence $\tau_i(q')=Y_i$ for $i\le n$,
$\rest^n(q')=q$, and $\tau_{n+i}(q')=\tau_i(q)\conv Y_{n+i}$ for $i\ge1$. By Theorem~\ref{thm:main},
$q'\conv\pK$, and Lemma~\ref{lem:inv}(b) gives $q=\rest^n(q')\conv\rest^n(\pK)$.
\end{proof}

\section{Closing}

Whether every run without an hnf has its teeth's reducts inside a set with passive argument
slots, like $\B$, is open. By
\citep[Theorem~7 and Proposition~9(b)]{s3}, two terms without an hnf whose teeth are pairwise
convertible have the same infinitary normal form, so CT is a case of the conjecture of~\citep{s6} that
$t^\infty=u^\infty$ implies $t\conv u$.

\begin{conjecture}
If $p,q\in\Ter$ have no hnf and $\tau_n(p)\conv\tau_n(q)$ for all $n\ge1$, then $p\conv q$.
\end{conjecture}

\section*{Data availability}
\sloppy
The proof is not mechanized. Its inputs are proved in the Lean~4~\citep{lean4} development
\texttt{lean-ski}~\citep{leanski}, version 1.20.0, for $SKI$-terms, with only the axioms
\texttt{propext} and \texttt{Quot.sound}: (CR) and (Std) as \texttt{church\_rosser} and
\texttt{wsteps\_split} in \texttt{Ski/Weak.lean}, (F2) as \texttt{Full.run\_total} in
\texttt{Ski/Full.lean}, through the block model of~\citep{s1}. Independent programs that follow
random reductions from $\pK$ and its teeth found every term reached in $\Red(\B)$; the proof does
not depend on them.

\bibliographystyle{plainnat}
\bibliography{refs}
\end{document}
